% !TEX program = xelatex
\documentclass[UTF8,zihao=5,a4paper,fontset=none]{ctexart}

\usepackage[
  left=1.2cm,
  right=1.2cm,
  top=0.9cm,
  bottom=0.9cm,
]{geometry}
\usepackage{array}
\usepackage{enumitem}
\usepackage{graphicx}
\usepackage{fontawesome5}
\usepackage{hyperref}
\usepackage{tabularx}
\usepackage{tikz}
\usepackage{xcolor}
\usepackage{xeCJKfntef}

\hypersetup{
  colorlinks=true,
  urlcolor=black,
  linkcolor=black,
}

% 优先使用跨平台字体，再回退到 macOS 常见字体或 TeX Live 自带的 Fandol 字体。
\IfFontExistsTF{TeX Gyre Termes}{\setmainfont{TeX Gyre Termes}}{\setmainfont{Times New Roman}}
\IfFontExistsTF{Noto Serif CJK SC}{%
  \setCJKmainfont{Noto Serif CJK SC}
}{%
  \IfFontExistsTF{Songti SC}{\setCJKmainfont{Songti SC}}{\setCJKmainfont{FandolSong-Regular.otf}}
}
\IfFontExistsTF{Noto Sans CJK SC}{%
  \setCJKsansfont{Noto Sans CJK SC}
  \newCJKfontfamily\resumeheiti{Noto Sans CJK SC}
}{%
  \IfFontExistsTF{Heiti SC}{%
    \setCJKsansfont{Heiti SC}
    \newCJKfontfamily\resumeheiti{Heiti SC}
  }{%
    \setCJKsansfont{FandolHei-Regular.otf}
    \newCJKfontfamily\resumeheiti{FandolHei-Regular.otf}
  }
}

\linespread{1.15}

\pagestyle{empty}
\setlength{\parindent}{0pt}
\setlength{\parskip}{0pt}
\setlength{\tabcolsep}{0.35em}
\hyphenpenalty=10000
\exhyphenpenalty=10000
\xeCJKsetup{CJKecglue={\hskip 0.14em plus 0.03em minus 0.01em}}

% 模板颜色集中放在这里，后续只需要改这一处即可整体换风格。
\definecolor{resumeText}{HTML}{1E1E1E}
\definecolor{resumeMuted}{HTML}{555555}
\definecolor{resumeHeaderMuted}{HTML}{666666}
\definecolor{resumeRule}{HTML}{000000}
\definecolor{resumePhotoBg}{HTML}{F7F7F7}
\color{resumeText}

\newcommand{\resumeSection}[1]{%
  \vspace{0.45em}%
  {\resumeheiti\zihao{4}\bfseries\color{black}#1}\par
  \vspace{0.08em}%
  {\color{resumeRule}\hrule height 1.4pt}%
  \vspace{0.28em}%
}

\newcommand{\resumeContactRow}[4]{%
  #1\hspace{0.35em}#2 & #3\hspace{0.35em}#4\\[-0.08em]
}

\newcommand{\resumeContactIcon}[1]{%
  \vphantom{Tel:}{\color{resumeHeaderMuted}\raisebox{-0.08ex}[0pt][0pt]{\fontsize{10}{10}\selectfont #1}}%
}

\newcommand{\resumeUrl}[2]{%
  \href{#1}{\CJKunderline{#2}}%
}

\newcommand{\resumeEducation}[3]{%
  {\resumeheiti\zihao{-4}\bfseries #1} & #2 & \hfill #3\\
}

\newcommand{\experienceTitle}[3]{%
  \begin{tabularx}{\textwidth}{@{}>{\raggedright\arraybackslash}X r@{}}
    {\resumeheiti\zihao{-4}\bfseries #1}{\zihao{5}\color{resumeMuted}｜#2} &
    {\zihao{5}\bfseries\color{black}#3}\\
  \end{tabularx}\par
  \vspace{0.06em}%
}

\newcommand{\projectTitle}[5]{%
  \begin{tabularx}{\textwidth}{@{}>{\raggedright\arraybackslash}X r@{}}
    {\zihao{-4}\textbf{#1}}{\zihao{5}\color{resumeMuted}｜#2} & {\zihao{5}\bfseries\color{black}#3}\\[-0.05em]
    \if\relax\detokenize{#4}\relax
    \else
      \multicolumn{2}{@{}l@{}}{{\color{resumeMuted}项目链接: }\resumeUrl{#4}{#5}}\\
    \fi
  \end{tabularx}\par
  \vspace{0.06em}%
}

% 可选：在项目/开源条目后追加仓库 star 数，如 \projectStars{100+ stars}。
\newcommand{\projectStars}[1]{%
  {\zihao{5}\mdseries\color{resumeMuted}（#1）}%
}

\newcommand{\resumeLabel}[1]{%
  {\resumeheiti\bfseries #1：}%
}

\newcommand{\resumeItem}[2]{%
  \item {\resumeheiti\bfseries #1：}#2%
}

\newenvironment{resumeItemize}{%
  \begin{itemize}[
    label=\textbullet,
    leftmargin=1.2em,
    itemsep=0em,
    topsep=0em,
    parsep=0.20em,
    partopsep=0pt
  ]%
  \setlength{\emergencystretch}{0.5em}%
}{%
  \end{itemize}%
}

% 如需使用证件照，把图片保存为 assets/photo.png；文件不存在时自动显示占位框。
\newcommand{\resumePhotoFile}{assets/photo.png}
\newcommand{\resumePhoto}{%
  \IfFileExists{\resumePhotoFile}{%
    \includegraphics[height=2.5cm]{\resumePhotoFile}%
  }{%
    \fcolorbox{resumeRule}{resumePhotoBg}{%
      \parbox[c][2.45cm][c]{1.9cm}{\centering\small 证件照}%
    }%
  }%
}

\begin{document}

\begin{tikzpicture}[remember picture, overlay]
  \node[anchor=north east, inner sep=0pt]
    at ([xshift=-1.2cm,yshift=-0.9cm]current page.north east)
    {\resumePhoto};
\end{tikzpicture}

\begin{minipage}[c]{0.80\textwidth}
  {\zihao{-2}\textbf{你的姓名}}\par
  \vspace{0.22em}
  \begin{tabular*}{\textwidth}{@{}p{0.38\textwidth}@{\extracolsep{\fill}}p{0.56\textwidth}@{}}
    \resumeContactRow{\resumeContactIcon{\faPhone}}{000-0000-0000}{\resumeContactIcon{\faEnvelope}}{\resumeUrl{mailto:name@example.com}{name@example.com}}
    \resumeContactRow{\resumeContactIcon{\faGlobe}}{\resumeUrl{https://example.com}{example.com}}{\resumeContactIcon{\faGithub}}{\resumeUrl{https://github.com/yourname}{github.com/yourname}}
    \multicolumn{2}{@{}l@{}}{{\resumeheiti\bfseries\color{resumeHeaderMuted}目标岗位｜每周可到岗 X 天及以上｜到岗时间｜可连续实习 X 个月}}\\
  \end{tabular*}
\end{minipage}

\resumeSection{教育经历}
\begin{tabularx}{\textwidth}{@{}>{\raggedright\arraybackslash}p{6.0cm}>{\raggedright\arraybackslash}X>{\raggedleft\arraybackslash}p{3.0cm}@{}}
  \resumeEducation{示例大学{\color{resumeMuted}（985/双一流）}}{硕士｜人工智能学院｜人工智能专业}{\textbf{2023.9} - \textbf{2026.7}}
  \multicolumn{3}{@{}l@{}}{{\hspace{0.45em}\zihao{5}\color{resumeMuted}排名前 \textbf{5\%}、国家奖学金、预计毕业时间 2026 年 7 月}}\\
  \resumeEducation{示例大学{\color{resumeMuted}（双一流）}}{本科｜计算机学院｜软件工程专业}{\textbf{2019.9} - \textbf{2023.7}}
  \multicolumn{3}{@{}l@{}}{{\hspace{0.45em}\zihao{5}\color{resumeMuted}排名前 \textbf{10\%}、校级一等奖学金、优秀毕业生}}\\
\end{tabularx}

\resumeSection{技能}
\resumeLabel{编程}熟悉 Python、PyTorch，了解 PostgreSQL、Redis，能实现数据预处理、训练/微调与推理评测流程。\par
\resumeLabel{专业方向}熟悉 RAG 检索链路、Agent 工具调用、多轮记忆、LoRA 微调、vLLM 推理部署与自动评测闭环。\par
\resumeLabel{工程化}熟悉 FastAPI、Docker、Git、Linux 及 AI 编程工具，具备镜像构建、多服务编排、环境配置与问题排查能力。

\resumeSection{实习经历}
\experienceTitle{示例科技有限公司}{AI 应用开发实习生}{2025.6 - 2025.9}
\resumeLabel{示例产品一｜智能审核平台}\par
\begin{resumeItemize}
  \resumeItem{多模态审核链路}{负责业务文档的多模态审核链路，结合多模态大模型、目标检测与 OCR 复核生成可定位、可回溯的结构化报告，将典型单文档审核耗时由约 \textbf{2 小时}压缩至 \textbf{20 分钟以内}。}
  \resumeItem{长文档解析}{基于 MinerU、PyMuPDF 解析 PDF、Word 及在线文档，将章节、表格、图片与位置信息绑定；按章节分块并行审核长文档，并通过页码、坐标和 block\_id 回溯结果。}
  \resumeItem{Agent 与部署}{基于 LangGraph 构建多轮审核 Agent，通过 asyncio 与 WebSocket 实现批量并发和进度推送；使用 Docker Compose 编排 FastAPI、MySQL、向量库与对象存储等服务。}
\end{resumeItemize}
\vspace{0.12em}
\resumeLabel{示例产品二｜智能体开发平台}\par
\begin{resumeItemize}
  \resumeItem{RAG 检索链路}{负责多格式文档解析、切分、Embedding、向量化、Dense/BM25 混合检索、RRF 融合与 Rerank 精排，在内部评测集上将 Top-5 召回率由约 \textbf{72\%} 提升至 \textbf{94\%}。}
  \resumeItem{工作流运行时}{负责可视化 Agent 工作流与后端执行语义对齐，覆盖条件分支、循环、并行、表单中断恢复及工具调用等节点，维护上下文、运行记录与错误传播。}
  \resumeItem{应用发布与扩展}{实现草稿调试与不可变版本快照，通过 Provider 插件统一接入 Chat、Embedding 与 Rerank 模型，实现 MCP/Skill、OpenAI 兼容接口和 SSE 流式输出；完善 Dockerfile 与锁文件，维护稳定部署流程。}
\end{resumeItemize}

\resumeSection{项目经历}
\projectTitle{多轮咨询的智能问答助手}{个人全栈开发}{2025.1 - 2025.3}{https://github.com/yourname/project}{github.com/yourname/project}
\resumeLabel{项目简介}面向多轮次、信息不完整的咨询场景，构建“信息采集 + 业务 Skill + 专业 Agent + 协作编排”的问答系统，解决背景信息缺失、复杂问题拆解不足和安全边界不稳定的问题。\par
\vspace{0.12em}
\begin{resumeItemize}
  \resumeItem{分层架构}{将检索、评估、建议和研究等能力拆分为 \textbf{9} 个业务 Skill；定义路由 Agent 和 \textbf{3} 个 Worker Agent，实现路由拆解与能力调用的职责解耦。}
  \resumeItem{路由机制}{构建“信息采集 → 快速咨询 → 协作编排”的三路路由机制：背景不足时由 Agent 逐轮追问关键信息；常见低风险咨询进入单 Agent 快速路径，复杂问题由主 Agent 拆解后交给多个 Worker Agent 并行处理；路由准确率达 \textbf{95\%}，并相对基线将常见咨询延迟成本降低约 \textbf{75\%}。}
  \resumeItem{记忆与安全}{实现短期会话历史与长期记忆机制，维护会话内最近 \textbf{5} 轮关键上下文，支持跨会话相似案例检索；通过 Harness 约束 Agent 能力边界、输出格式与安全策略，并自动修复不合规输出，综合得分达到 \textbf{4.5/5}。}
  \resumeItem{评测闭环}{在自建的多轮咨询样例集上结合 LLM-as-a-Judge 与人工抽检评估，相比基线显著提升最终回答的完整性与稳定性，并沉淀可复现的评测脚本与样例集。}
\end{resumeItemize}

\resumeSection{开源贡献}
\projectTitle{开源项目示例一\hspace{0.3em}\projectStars{100+ stars}}{Merged PR Contributor}{2025.4}{}{}
\resumeLabel{贡献}(\resumeUrl{https://github.com/owner/repo/pull/1}{\#1}) 修复终端未加载用户 shell 配置的问题，统一使用 POSIX \texttt{sh} 执行脚本，补充回归测试。\par
\vspace{0.12em}
\projectTitle{开源项目示例二\hspace{0.3em}\projectStars{50+ stars}}{Merged PR Contributor}{2025.5 - 2025.6}{}{}
\resumeLabel{贡献}(\resumeUrl{https://github.com/owner/repo/pull/2}{\#2}、\resumeUrl{https://github.com/owner/repo/pull/3}{\#3}) 修复生产环境可选依赖安装；为内置子 Agent 增加隔离上下文并兼容严格模型提供商，补充回归测试。\par

\end{document}
